\chapter{Appendix B: Running, Testing, and Hacking}

This appendix is a practical companion.
It explains how to work with the repository as you read.

\section{Where the code lives}
The canonical repository is:
\href{https://github.com/ravikisha/relax.js}{https://github.com/ravikisha/relax.js}.

The published NPM package is named \texttt{relaxcore}:
\href{https://www.npmjs.com/package/relaxcore}{https://www.npmjs.com/package/relaxcore}.

If you are reading the book to learn, I recommend working from source.
If you are integrating Relax.js into another project, use the NPM package.

\section{The contract: scripts are interfaces}
In this repo, scripts in \texttt{package.json} are not convenience shortcuts.
They are stable interfaces for humans and for automation.

Typical tasks you\textquotesingle ll do while reading:

\begin{itemize}
  \item Run unit tests for a chapter and then break one behavior on purpose.
  \item Run lint/typecheck to keep refactors honest.
  \item Build ESM/CJS outputs and verify the exports map.
\end{itemize}

\section{Suggested local workflow}
Treat the following loop as your default while experimenting:

\begin{enumerate}
  \item Run the full test suite.
  \item Run only the relevant test file for the chapter you\textquotesingle re reading.
  \item Make one small change.
  \item Re-run the focused test first, then the full suite.
\end{enumerate}

\subsection{A minimal command cheat sheet}
These commands are intentionally boring.
They keep you honest.

\begin{itemize}
  \item Install: \texttt{npm install}
  \item Run all tests: \texttt{npm test}
  \item Run tests in watch mode (if configured): \texttt{npm run test:watch}
  \item Lint: \texttt{npm run lint}
  \item Build: \texttt{npm run build}
\end{itemize}

If a command name differs in your local setup, use \texttt{package.json} as the source of truth.

\section{Debugging strategy: shrink the mystery}
When something fails, don\textquotesingle t guess.
Reduce the problem:

\begin{enumerate}
  \item Find the smallest test that demonstrates the failure.
  \item Add logging or instrumentation (then remove it).
  \item Make the smallest change that fixes the bug.
  \item Add a regression test.
\end{enumerate}

\section{Common failure modes}
\subsection{DOM tests behave differently than browser demos}
A jsdom environment is not a full browser.
When a behavior is browser-specific, keep the boundary explicit.

\subsection{Performance claims feel inconsistent}
Performance depends on machine, browser version, and workload.
Use the benchmark methodology chapter and prefer relative comparisons within the same environment.

\subsection{Packaging and exports feel confusing}
Modern Node resolution uses \texttt{exports} maps.
If an import fails, verify the entry point exists in the built output and that it is mapped correctly.
See: \href{https://nodejs.org/api/packages.html\#exports}{https://nodejs.org/api/packages.html\#exports}.
